\documentclass[UTF8,fontset=none,11pt,a4paper]{ctexart}
\usepackage{ipara-handbook}
\handbooksetup{DS-04 树与二叉树}{408 · 数据结构 DS-04}{Hierarchy · Traversal · Ownership · Coding}
\begin{document}
\handbooktitle{树与二叉树}{层次关系怎样递归表示、遍历并保持结构不变量}{408 · 数据结构 Topic DS-04}

\begin{corebox}{母问题：层次关系怎样变成可执行的访问顺序？}
树的生成链是
\[
\text{Hierarchy}\to\text{Recursive Representation}\to\text{Frontier Traversal}\to\text{Local Change}\to\text{Structural Invariant}.
\]
节点的子树仍是树，所以递归定义自然产生递归算法；把尚未展开的节点放入栈或队列，就得到显式遍历。访问时机决定先序、中序、后序，取出规则决定层序。Huffman 则把“树”用于编码优化：反复合并最小权值，目标是最小带权路径长度。
\end{corebox}
\begin{methodbox}{本册定位与 Owner 边界}
\begin{itemize}
\item \kw{Owns：}树/森林与二叉树的性质和存储、孩子兄弟转换、递归与迭代遍历、线索化、高度、Huffman 主干。
\item \kw{Uses：}DS02 的栈/队列端点契约；DS05 的优先队列接口。
\item \kw{Stops：}堆内部上滤下滤、BST/AVL/B/B+ 查找与旋转、图遍历正确性归各自 Topic。
\end{itemize}
\end{methodbox}
\tableofcontents\newpage

\section{对象与表示}
二叉树节点至少包含值、左子树、右子树。子树为空是合法状态，不是异常；叶节点的两个孩子都为空。用唯一所有权表示时，父节点拥有子节点，销毁根即可递归释放整棵树。树的结构不变量是：每个节点至多有一个父关系入口，沿孩子指针不会回到祖先，且遍历所得节点数等于可达节点数。

\begin{examplebox}{同一棵树的四种读法}
对根 1、左子树 2(4,5)、右子树 3，先序为 $1,2,4,5,3$，中序为 $4,2,5,1,3$，后序为 $4,5,2,3,1$，层序为 $1,2,3,4,5$。树没有变化，改变的是“何时从 frontier 取出并处理根”。
\end{examplebox}

\section{遍历：访问时机与 frontier}
递归先序在进入节点时处理；中序在左子树返回后处理；后序在两个子树都返回后处理。显式栈保存尚未完成的调用现场：先序压入右子树再压左子树，才能保证左子树先出栈。中序先沿左链压栈，弹出一个节点后转向右子树。层序使用 FIFO 队列，按深度展开；其队列语义来自 DS02，树的覆盖正确性由本册维护。

\section{树、森林与二叉树：孩子兄弟表示是可逆接口}
一般树的节点可以有任意多个孩子。孩子兄弟表示只保留两个方向：
\[
\code{firstChild}=\text{第一个孩子},\qquad
\code{nextSibling}=\text{下一个兄弟}.
\]
把 \code{firstChild} 解释为二叉树左孩子、\code{nextSibling} 解释为右孩子，就得到树到二叉树的转换。森林则把各棵树的根按兄弟链连接：第一棵树根是二叉树根，后续树根依次位于右链。

该转换保持层次与兄弟次序，不保持“二叉树左右孩子”的原语义。由指针解释可推出遍历对应：
\[
\text{树的先根遍历}=\text{对应二叉树的先序遍历},\qquad
\text{树的后根遍历}=\text{对应二叉树的中序遍历}.
\]
森林遍历还要沿根的兄弟链完成每棵树，不能只处理第一棵。

\section{线索二叉树：把空指针改造成遍历序列邻接关系}
含 $n$ 个节点的二叉链表有 $2n$ 个孩子指针，而非空边只有 $n-1$ 条，因此有 $n+1$ 个空指针。线索化利用这些空槽保存某种遍历次序下的前驱或后继。因为字段可能表示真实孩子或线索，必须增加标志：
\[
ltag=0\Rightarrow lchild\text{ 是左孩子},\quad
ltag=1\Rightarrow lchild\text{ 是前驱线索},
\]
右侧同理。忽略 tag 沿线索递归，会把序列关系误当结构边，造成重复访问甚至循环。

\subsection{中序线索化的状态机}
维护 $pre$ 指向刚刚访问的节点，在访问当前节点 $p$ 时：
\begin{enumerate}
\item 若 $p$ 的左孩子为空，令 $p.lchild=pre$ 且 $p.ltag=1$；
\item 若 $pre$ 存在且其右孩子为空，令 $pre.rchild=p$ 且 $pre.rtag=1$；
\item 更新 $pre=p$，再处理右子树。
\end{enumerate}
这不是任意补指针，而是在同步构造中序序列相邻节点的关系。线索化完成后找中序后继：若 $rtag=1$，线索直接给出后继；若 $rtag=0$，后继是右子树中沿真实左孩子到达的最左节点。

\begin{warnbox}{线索树边界}
线索只服务于指定遍历次序；中序前驱/后继不能直接当作先序关系。某些先序、后序邻接仍可能需要父指针或重新定位，不能把“有线索”理解成所有遍历操作都严格 $O(1)$。
\end{warnbox}

\section{Huffman：树是合并结果，不是搜索树}
给定符号权重，带权路径长度为 $\sum_i w_i d_i$。每次合并两个最小权值 $x,y$，把 $x+y$ 重新放回候选集合，累计代价 $x+y$。该贪心过程依赖最小优先队列；本册只固定合并契约与代价，优先队列内部实现由 DS05 Own。单个符号的编码长度约定为 0，空输入代价为 0。

\section{代码：从节点所有权到访问顺序}
完整实现位于 \codepath{code/ds04_tree_binary.hpp}，测试位于 \codepath{code/ds04_tree_binary_test.cpp}；代码块直接引用真实源文件。
\subsection{节点与递归遍历}
\lstinputlisting[language=C++,firstline=12,lastline=43,firstnumber=12,numbers=left,caption={唯一所有权节点与递归先中后序}]{30_408/10_数据结构/04_树与二叉树/code/ds04_tree_binary.hpp}
递归基是空指针立即返回；每个非空节点只处理一次，子树调用返回后再按访问时机写入输出。
\subsection{显式栈与层序队列}
\lstinputlisting[language=C++,firstline=45,lastline=99,firstnumber=45,numbers=left,caption={迭代遍历与层序 frontier}]{30_408/10_数据结构/04_树与二叉树/code/ds04_tree_binary.hpp}
显式结构保存的不是“所有节点”，而是尚未展开的边界；栈的压入顺序和队列的端点顺序直接生成输出次序。
\subsection{Huffman 合并调用}
\lstinputlisting[language=C++,firstline=111,lastline=130,firstnumber=111,numbers=left,caption={最小权值合并与累计代价}]{30_408/10_数据结构/04_树与二叉树/code/ds04_tree_binary.hpp}
\subsection{测试证据}
\lstinputlisting[language=C++,firstline=17,lastline=52,firstnumber=17,numbers=left,caption={空树、遍历覆盖与 Huffman 边界测试}]{30_408/10_数据结构/04_树与二叉树/code/ds04_tree_binary_test.cpp}
\begin{lstlisting}[language=bash]
clang++ -std=c++17 -Wall -Wextra -Werror -pedantic \\
  code/ds04_tree_binary_test.cpp -o /tmp/ds04_test
/tmp/ds04_test
\end{lstlisting}

\section{复杂度与概念边界}
\begin{center}\small
\begin{tabularx}{\linewidth}{L{2.5cm}C{1.8cm}C{1.8cm}YY}\toprule
操作 & 递归/迭代遍历 & 额外空间 & 前提 & 边界\\\midrule
遍历 $n$ 节点 & $O(n)$ & $O(h)$ 或 $O(w)$ & 每节点一次 & 空树返回空序列\\
高度 & $O(n)$ & $O(h)$ & 子树高度定义固定 & 空树高度约定\\
Huffman 合并 & $O(k\log k)$ & $O(k)$ & 最小优先队列 & $k<2$ 代价为 0\\\bottomrule
\end{tabularx}\end{center}
\begin{center}\small
\begin{tabularx}{\linewidth}{L{1.9cm}C{.35cm}L{1.9cm}YY}\toprule
概念 A&$\neq$&概念 B&判据&错误后果\\\midrule
树节点所有权&$\neq$&遍历顺序&生命周期 vs 访问时机&泄漏或次序错\\
层序遍历&$\neq$&BFS 图遍历&树无环且无需一般 visited&把图结论硬套树\\
Huffman 树&$\neq$&BST&最小带权路径 vs 有序查找&混淆目标\\
\bottomrule\end{tabularx}\end{center}

\section{做题控制协议与压缩复原}
先画根、子树范围和空孩子；再选访问时机或 frontier；每次节点只入队/入栈一次并验证覆盖；若题目问 Huffman，改写为“合并两个最小权值并累计”。忘掉口诀后复原五问：\kw{节点拥有谁？未展开边界在哪里？当前节点何时处理？空子树如何结束？目标是访问次序还是编码代价？}

\begin{corebox}{全科交接}
DS04 从 DS02 接收 LIFO/FIFO frontier，从 DS05 接收最小优先队列接口；向 DS-B01 提供树遍历实例，向 DS09 提供“树结构可承载有序索引”的前置直觉，但不拥有 BST/AVL 机制。
\end{corebox}
旧总册关于递归层次、遍历访问时机与 Huffman 合并的内容已吸收；树/森林转换与线索化已按考纲补入本册，遍历序列重建和更复杂编码变体保留为扩展，不能覆盖本册主干。
\end{document}
